\section{Fermionic sphere tensors}\label{sec:fermions}

We need a tensor whose legs can be combined when their positions
coincide. Start with marked points $z_1,\ldots,z_k$ in an affine
coordinate on the projective line. At each point the input space has a
half-integral \emph{energy} $h_i$. The energies fix the dimensions of
the legs, while the marked points determine the entries of the
tensor. The metric sites at which we will price these legs can be
chosen independently of the $z_i$.

Two features of this construction enter the proof. When a cluster of
marked points collides, the Laurent coefficient of a given order
passes through a single leg, whose energy the order determines when the other sites stay fixed
(Proposition~\ref{prop:fusion}).
Conversely, at a two-point collision, varying the first input across
a bounded range of energies realizes every linear map from the second
input to the fused leg. Lemma~\ref{lem:width} gives that energy bound.
Together these facts let Section~\ref{sec:sphere} compare the price
of a sphere tensor with the price of moving an independent wire.

The half-integral grading comes from the free fermion
\cite[Section~4]{kacwang}. Khoruzhii used its modes to decompose a
crossing of two wires \cite{Khoruzhii}. The crossing argument
of \cite[Theorem~\ext{G}{thm:crossing}]{oaG} uses the same calculus
in arbitrary characteristic.

\subsection{The kernel and its coefficients}
Let $F$ be a field and take the exterior algebra on generators $\psi_m$,
$m=0,1,\ldots$. Give $\psi_m$ energy $m+\tfrac12$, and denote by $V_h$
the subspace of energy $h\in\frac12\mathbb Z_{\ge0}$. Its basis consists of
the products $\psi_A=\bigwedge_{m\in A}\psi_m$, in increasing order, with
$\sum_{m\in A}(m+\tfrac12)=h$. In particular $V_0=F$. Every such product
has exterior degree congruent to $2h$ modulo two.

For distinct $z_1,\ldots,z_k$ in $F$ we define a multilinear form
$C(\mathbf h;\mathbf z)$ on $\prod_i V_{h_i}$. For basis inputs, list all generators in site order.
The value
is a signed sum over perfect matchings of this list (Wick's
rule~\cite{wick}). Its factors are as follows. Two generators at one
site have contraction zero. At different sites, the contraction of
$\psi_m$ at $z_i$ with $\psi_n$ at $z_j$, in that order, is
\begin{equation}\label{eq:fermions-kernel}
 [x^m y^n]\,\frac{1}{z_j+y-z_i-x}
 =\frac{(-1)^n\binom{m+n}{m}}{(z_j-z_i)^{m+n+1}} .
\end{equation}
Pair the entries of each oriented matching consecutively, multiply
their contractions, and insert the sign of the permutation from
the original list to this paired list. Reversing a pair reverses both
the contraction and the permutation sign. Changing the order of pairs
preserves the sign. This sum is the Pfaffian of the alternating
matrix of contractions. The empty list has value one and an odd
list has value zero.

For example, put one generator $\psi_0$ at each of four sites
and write $K_{ij}=(z_j-z_i)^{-1}$. The value on these four inputs is
\[
 K_{12}K_{34}-K_{13}K_{24}+K_{14}K_{23}.
\]
At a collision of $z_1$ with $z_2$ the first matching has an internal
pair and a pole, while the other two matchings connect the cluster to
the outside. This distinction governs the general fusion expansion.

These formulas use integer coefficients and commute with change of field.
Two repeated generators at a site give canceling pairs of matchings, by
exchanging their partners. Thus the forms descend to exterior powers in
characteristic two as well. Interchanging sites $i,j$ in adjacent positions
contributes the sign $(-1)^{(2h_i)(2h_j)}$. Each coefficient, viewed as a rational function,
is homogeneous of degree $-\sum_i h_i$ in the positions, since a contraction
has degree minus the sum of its two energies. Its only possible poles are
site collisions. As a function of $z_i$ it is bounded at infinity, and
vanishes there when $h_i>0$.

\subsection{Fusing a cluster at a marked point}
Partition the sites into consecutive clusters $J_1,\ldots,J_s$, reordering
the sites first if necessary. In cluster $J_j$ write $z_i=p_j+u a_i$, where
the centers $p_j$ are distinct and the offsets within each cluster are
distinct. A singleton with offset zero stays fixed. Its fusion map is the
identity at its original energy and zero at every other energy.

\begin{proposition}[Fusion]\label{prop:fusion}
For each cluster $J$ there are maps
\[
 I_e^J:\bigotimes_{i\in J}V_{h_i}\longrightarrow V_e
 \qquad(e\in\tfrac12\mathbb Z_{\ge0})
\]
depending only on its offsets and input energies. Put
$H_J=\sum_{i\in J}h_i$. The maps vanish unless $e-H_J\in\mathbb Z$.
The Laurent expansion at $u=0$ is
\begin{equation}\label{eq:fermions-fusion}
 C(\mathbf h;\mathbf z)
 =\sum_{\mathbf e}u^{\,\sum_j e_j-\sum_i h_i}\,
 C(\mathbf e;\mathbf p)\circ
       (I^{J_1}_{e_1},\ldots,I^{J_s}_{e_s}).
\end{equation}
Only finitely many terms contribute to each coefficient.
\end{proposition}

\begin{proof}
Group the perfect matchings by the pairs they use inside each cluster.
Fix one cluster $J$ and one such partial matching. Contract its pairs
using \eqref{eq:fermions-kernel} at positions $a_i$. Move the matched pairs to the front, with their permutation sign,
keeping all unmatched generators in order. Replace each unmatched generator
at offset $a_i$ by
\begin{equation}\label{eq:fermions-translation}
 \psi_m\longmapsto
          \sum_{\ell\ge m}\binom{\ell}{m}a_i^{\ell-m}\psi_\ell .
\end{equation}
Wedge the resulting vectors in their original order and keep only energy
$e$. Each translation may be truncated to the modes of energy at most
$e$, so this is an operation on finite-dimensional spaces. Sum over
all internal matchings with their signed factors to obtain $I_e^J$.
If a generator is repeated in one input and either copy is contracted,
interchanging the copies pairs terms of opposite signs. If both are
unmatched, their translated wedge vanishes. This cancellation
works over every field. Removing pairs
and raising mode indices preserve energy modulo $\mathbb Z$.

Decompose each full matching into its internal edges and its edges between
clusters. An internal contraction at the original positions contributes the
stated contraction at the offsets times $u$ raised to minus the energy of
the pair. On an edge between clusters the kernel is regular at $u=0$.
Taylor expansion in the two offsets uses
$[x^m](u a_i+x)^\ell=\binom{\ell}{m}(u a_i)^{\ell-m}$, which gives
\eqref{eq:fermions-translation} and one factor of $u$ per index increase.
For a given output energy $e$, the total power from the cluster is therefore
$e-H_J$. Moving the internal pairs of all clusters to the front contributes
only the signs already included in $I_e^J$, since an ordered pair can
cross any earlier cluster without changing sign. What remains is the
sign of the matching between the output wedges. This proves
\eqref{eq:fermions-fusion}. At fixed total power the sum of the nonnegative
$e_j$ is fixed, so the coefficient is finite.
\end{proof}

For example, fuse two generators $\psi_0,\psi_0$ at offsets $0,1$.
Their total energy is $1$. The internal contraction produces the vacuum
at energy $0$ with coefficient $u^{-1}$. If the two generators remain
unmatched, their first nonzero translated wedge is
$\psi_0\wedge\psi_1$, of energy $2$ and order $u^1$. This second term
can pair with generators at other centers.

\subsection{Energy and dimension}
A leg of dimension $d_h=\dim V_h$ carries width $\log d_h$ when
$d_h>0$. A basis wedge chooses a subset of distinct modes, and its
energy is the sum of their half-integer weights. Counting these
subsets gives width asymptotic to $c\sqrt h$, where $c=\pi/\sqrt3$
and the logarithm is natural.

\begin{lemma}[Dimension of a level]\label{lem:dimension}\label{lem:fermions-dimension}
For every $h\in\frac12\mathbb Z_{\ge0}$ one has $d_h\le e^{c\sqrt h}$.
As $h\to\infty$ through all half-integral levels,
\[
   d_h\sim
   \frac{\exp(c\sqrt h)}{2^{5/2}3^{1/4}h^{3/4}} .
\]
\end{lemma}

\begin{proof}
For $t>0$ the generating function is
\[
 Z(t)=\sum_h d_h e^{-th}
     =\prod_{m\ge0}\bigl(1+e^{-t(m+1/2)}\bigr).
\]
The function $f(x)=\log(1+e^{-x})$ is convex on $[0,\infty)$.
The midpoint estimate on each interval $[mt,(m+1)t]$ gives
$\log Z(t)\le b/t$, with
$b=\int_0^\infty f(x)\,dx=\sum_{j\ge1}(-1)^{j+1}/j^2=\pi^2/12$.
Thus $d_h\le \exp(th+b/t)$. Minimization in $t$ gives the bound for $h>0$,
and $d_0=1$.

Doubling the generator energies identifies basis wedges at energy $h$ with
partitions of $2h$ into distinct odd parts. These are in bijection with the
self-conjugate Young diagrams of size $2h$. Each diagonal hook has length
$2m+1$, where $m$ is its arm length, and the strictly decreasing sequence of
arm lengths reconstructs the diagram. The asymptotic formula for the number
of self-conjugate partitions of $n$ is
$\exp(\pi\sqrt{n/6})/(2^{7/4}3^{1/4}n^{3/4})$ to leading
order~\cite[Section~5.1]{Wildon}. Set $n=2h$.
\end{proof}

\subsection{The width needed to transmit a leg}
Consider a two-site cluster at offsets $0,1$ with incoming energies
$h,g$. Write $I_e$ for its fusion map to energy $e$. Fixing the first
input turns this map into an operator on the second input. To recover
a wire in Section~\ref{sec:sphere}, we need every matrix unit from
$V_g$ to $V_e$. The next lemma constructs them with first-input width
at most $c(\sqrt e+\sqrt g)+O(1)$ per level.

\begin{lemma}[Fusion width]\label{lem:width}
For $e,g\in\frac12\mathbb Z_{\ge0}$ put
\[
 H(e,g)=(\sqrt e+\sqrt g)^2+(\sqrt e+\sqrt g)/\sqrt2.
\]
The map
\[
 \bigoplus_{h\le H(e,g)} V_h \longrightarrow
       \operatorname{Hom}(V_g,V_e),\qquad
 (v_h)_h\longmapsto\sum_h I_e(v_h,-)
\]
is surjective. In the wedge bases, each matrix unit has an integral preimage.
\end{lemma}

\begin{proof}
Fix mode sets $A,B$ of energies $e,g$, respectively. We construct
an input at the first point that sends $\psi_B$ to $\psi_A$ and
annihilates every other basis vector of $V_g$. If $e=g=0$, the first
input is the vacuum. Otherwise write $w=\sqrt e$ and $v=\sqrt g$.
An input at offset zero will be a wedge of test vectors. Represent a test
$\sum_m t_m\psi_m$ by its Laurent polynomial $f(X)=\sum_m t_m X^{-m-1}$.
In a matrix coefficient of $I_e$, the test has two possibilities.
It can contract with input index $b$ at offset one, with coefficient
$[y^b]f(1+y)$ by \eqref{eq:fermions-kernel}. Or it can occupy output index
$\ell$, with coefficient $t_\ell$, because translation at offset zero is
the identity. This follows by expanding the output wedge. Each test must
occupy one slot of one of these two kinds, and different tests use different
slots.

We choose tests that can use their designated slots or slots of strictly
larger weight. The target slots will account for the entire available
weight, forcing every test into place. Give output index $\ell$ the
weight $\tau=(\ell+\tfrac12)/w$ and input index $b$ the weight
$\tau=(b+\tfrac12)/v$. A side of energy zero contains
no indices and requires no division. The test for each target slot is given
in the table.
\[
\begin{array}{c|c|c}
\text{target}& f(X)& r\\ \hline
\ell\in A&X^{-\ell-1}(1-X^{-1})^r&
   \lfloor v\tau-\tfrac12\rfloor+1\\
b\in B&X^{-r-1}(1-X^{-1})^b&
   \lfloor w\tau-\tfrac12\rfloor+1 .
\end{array}
\]
For an output target the coefficient $t_\ell$ is one, and the zero of order
$r$ at $X=1$ forces every input slot it can occupy to have larger weight.
Its other possible output indices exceed $\ell$. For an input target,
$f(1+y)=y^b(1+y)^{-r-b-1}$, so its leading input coefficient is one.
All other possible input slots have larger indices, and the choice of $r$
forces all possible output slots to have larger weight. These conclusions
hold over the integers.

Form the wedge of all the tests in a fixed order. In every matrix
coefficient of the indicated energies the sum of the weights of all
available slots is $w+v$, which is also the sum of the target weights.
The construction gives each test a slot of at least its target weight.
Since the slots have positive weights, every nonzero term must use
every slot and attain equality for each test. The table admits only
the designated slot at that weight, even if an opposite-side slot
has equal weight. In particular every generator of the second input
has contracted with a test. The input and output sets are therefore
$B$ and $A$, and the pairing and wedge assignment are unique. Thus
the desired coefficient is $1$ up to sign, which can be absorbed into the chosen input wedge.

The largest generator energy in each test of weight $\tau$ is at most
$(w+v)\tau+\tfrac12$, by the displayed choice of $r$. A wedge of $s$
generators has energy at least
$\sum_{m=0}^{s-1}(m+\tfrac12)=s^2/2$, so there are at most
$\sqrt2(w+v)$ tests. Adding their individual bounds shows that every
homogeneous component of the input wedge has energy at most $H(e,g)$.
Taking these components gives the claimed preimage.
\end{proof}

The sum of square roots in this cutoff also follows from a dimension count.
If the map in Lemma~\ref{lem:width} is surjective with cutoff $H$, then
$d_e d_g\le \sum_{h\le H}d_h\le (2H+1)e^{c\sqrt H}$.
If $H\le(\sqrt e+\sqrt g)^2$, taking logarithms and
using Lemma~\ref{lem:dimension} gives
$\sqrt H\ge\sqrt e+\sqrt g-O(\log(e+g+2))$ along nonzero levels
with $e+g\to\infty$. The same lower bound holds immediately for larger $H$.
Lemma~\ref{lem:width} attains this leading width in
every characteristic.
